\documentclass[twoside]{scrreprt}

%Packages, die für die deutsch Sprache erforderlich sind
\usepackage[utf8]{inputenc}
\usepackage[T1]{fontenc}
\usepackage{lmodern}
\usepackage{ngerman}

%Packages für Graphik
\usepackage[pdftex]{graphicx}
\graphicspath{{Abbildungen/} }

%Package, damit Bibtex-URL klappt
\usepackage{url}

\begin{document}
%%%%% BEGINN TITELSEITE %%%%%
\begin{titlepage}

% Oberer Teil der Titelseite:
\begin{minipage}{0.4\textwidth}
	\begin{flushleft}
		\includegraphics[height=1.5cm]{TU_Logo}
	\end{flushleft}
\end{minipage}
\hfill
\begin{minipage}{0.4\textwidth}
	\begin{flushright}
		\includegraphics[height=1.5cm]{MPM_Logo}%TODO: Logo in HQ
	\end{flushright}
\end{minipage}

\begin{center}    
\textbf{\LARGE Technische Universität Berlin}\\
\large{Institut für Konstruktion, Mikro- und Medizintechnik}\\
\large{Fachgebiet Methoden der Produktentwicklung und Mechatronik}\\
\large{\textsc{Prof. Dr.-Ing. Dietmar Göhlich}}\\
\vfill

% Title
\huge{\textbf{Vergleich von Ladesystemen und Speichertechnologien für batteriebetriebene Stadtbusse}}\\
\rule{\linewidth}{1pt}
\large{Zur Erlangung des akademischen Grades Bachelor of Science}\\
\large{Im Studiengang Informationstechnik im Maschinenwesen}\\
\vfill

%Autor
\large{\textsc{Ludger Heide}}\\
\large{Matrikelnummer 338641}\\
\vfill

\large{Betreuerin: \textsc{M.Sc. Tu-Anh Ly}}\\
\vfill

\large{\today}
\end{center}
\end{titlepage}
%%%%% ENDE TITELSEITE %%%%%

\pagenumbering{Roman}
\chapter*{Aufgabenstellung}

\chapter*{Eidesstattliche Erklärung}
Ich, Ludger Heide, versichere hiermit, dass ich meine Bachelorarbeit mit dem Thema
\begin{quote}
	\emph{Vergleich von Ladesystemen und Speichertechnologien für batteriebetriebene Stadtbusse}
\end{quote}
selbstständig verfasst und keine anderen als die angegebenen Quellen und Hilfsmittel benutzt habe, wobei ich alle wörtlichen und sinngemäßen Zitate als solche gekennzeichnet habe. Die Arbeit wurde bisher keiner anderen Prüfungsbehörde vorgelegt und auch nicht veröffentlicht.\\[6ex]
Berlin, \today\\
\newline
\rule{4cm}{0.5pt}\\
\textsc{Ludger Heide} 

\tableofcontents

\listoffigures

\listoftables
\newpage

%TODO: Abstract

%%%%% BEGINN INHALT %%%%%
\pagenumbering{arabic}
\chapter{Einleitung}
\section{Kontext} %Namen ändern!
\subsection{Geschichte}
\subsection{Aktueller Stand}
\section{Problem}
%TODO: WIchtigste Aufgabe sind die Ladesysteme
\subsection{Ziel} %TODO: Wegmachen?
%Neues Kapitel: Methodenteil?
\section{Vorgehen} %Methodik?
\section{Überblick} %TODO Wirklich? (lutz sagt nein)

\chapter{Ladesysteme} %Zewitens
\section{Bewertungskriterien} %TODO: Subsections
% Mit welcher Methode ausgewählt?
\section{Betrachtete Systeme} %TODO: Subsections
\section{Vergleichstabelle}   %TODO: In den Anhang?

\chapter{Speichertechnologien} %Zuerst
\section{Bewertungskriterien} %TODO: Subsections
\section{Betrachtete Systeme} %TODO: Subsections
\section{Vergleichstabelle}   %TODO: In den Anhang?

\chapter{Effizienzberechnung} %TODO: Besserer Titel %Drittens
\section{Ladesystem}
\section{Speichertechnologie}
\section{Ladestrategie und Route}
\section{Ergebnisse}
\subsection{Route A}
\subsection{Route B}

\chapter{Bewertung und Diskussion} %Diksussion
\section{Fazit}
\section{Ausblick}


%%%%% ENDE INHALT %%%%%

%Bibliographie
\bibliographystyle{alphadin}
\bibliography{Quellen/Quellenliste} 


\end{document}